%%%%%%%%%%%%%%%%%%%%%%%%%%%%%%%%%%%%%%%%%%%%%%%%%%%%%%%%%%%%%%%%%%%%%%%%%%%
%  ARCHIVED -- superseded by docs/chapter/Capacity_Is_Not_Access.tex.
%  Kept for the record only. Do not edit, compile for submission, or cite
%  numbers from this version; its results predate the corrected-graph runs.
%%%%%%%%%%%%%%%%%%%%%%%%%%%%%%%%%%%%%%%%%%%%%%%%%%%%%%%%%%%%%%%%%%%%%%%%%%%
%  2026 Summer Proceedings -- student chapter
%  "More Beds or Better Beds?"  Fatima Asghar
%  Built on the official svmult template. Compile with pdfLaTeX:
%      pdflatex chapter -> bibtex chapter -> pdflatex chapter -> pdflatex chapter
%%%%%%%%%%%%%%%%%%%%%%%%%%%%%%%%%%%%%%%%%%%%%%%%%%%%%%%%%%%%%%%%%%%%%%%%%%%

\documentclass[graybox]{svmult}

\usepackage{type1cm}          % activate if the standard fonts are unavailable
\usepackage{makeidx}          % allows index generation
\usepackage{graphicx}         % standard LaTeX graphics tool
\usepackage{multicol}         % used for the two-column index
\usepackage[bottom]{footmisc} % places footnotes at page bottom
\usepackage{url}              % the template recommends this for web addresses

\usepackage{newtxtext}        %
\usepackage{newtxmath}        % selects Times Roman as basic font

\graphicspath{{figures/}}

\makeindex

%%%%%%%%%%%%%%%%%%%%%%%%%%%%%%%%%%%%%%%%%%%%%%%%%%%%%%%%%%%%%%%%%%%%%%%%%%%

\begin{document}

\title*{More Beds or Better Beds? Agent-Based Modeling of Clean-Air Shelter
Access for Unsheltered Residents During Wildfire Smoke}
\titlerunning{More Beds or Better Beds?}
\author{Fatima Asghar}
\authorrunning{Fatima Asghar}
\institute{Fatima Asghar \at Harrisburg Area Community College, Harrisburg, PA,
\email{fxa28196@hawkmail.hacc.edu}}

\maketitle

\abstract{Wildfire smoke is now a recurring hazard in the Pacific Northwest, and
people living unsheltered cannot go indoors when it arrives. This chapter asks
which intervention helps them most: more clean-air shelter capacity, or
better-placed capacity. I built an agent-based model in Repast Simphony in which
6,842 unsheltered residents---the 2025 Point-in-Time count for Multnomah
County---start at real reported campsite locations, carry individually sampled
age, sex, mobility, asthma and COPD attributes, and walk Portland's real street
network toward the nearest shelter with room, driven by measured hourly
PM$_{2.5}$ from the September 2020 smoke event. Three scenarios were compared
across nine random seeds each. Today's 36-facility system shelters 30.1\% of
residents. Raising capacity to meet demand at the same locations shelters
91.6\%, but leaves 578 spaces empty while turning 562 people away. Spending that
identical capacity on ten optimally placed new sites shelters 96.0\% and halves
refusals. Capacity expansion alone widens the gap between residents who walk
easily and those who do not, from 13.0 to 24.5 percentage points; well-placed
capacity returns it to 12.9. No health outcome is modelled, and modelled uptake
is an upper bound.}

%%%%%%%%%%%%%%%%%%%%%%%%%%%%%%%%%%%%%%%%%%%%%%%%%%%%%%%%%%%%%%%%%%%%%%%%%%%
\section{Introduction}
\label{sec:intro}

In September 2020, wildfire smoke settled over Portland, Oregon for nine days.
Hourly PM$_{2.5}$ measured by the county's regulatory monitors peaked at
562.7~$\mu$g~m$^{-3}$ and stayed above the U.S.\ Environmental Protection
Agency's ``Unhealthy'' threshold for 194 of 312 hours. Oregon's Department of
Environmental Quality reports that this was Portland's first experience of air
that bad from wildfire: the city had recorded no unhealthy-or-worse smoke days
between 1985 and 2014, and in 2020 it recorded three very unhealthy days and
five hazardous ones \cite{deq2023trends}.

The public-health advice during such an episode is to stay indoors. That advice
is unavailable to people who have no indoors. Multnomah County opened clean-air
shelters in response, and the question this chapter addresses is what would
happen if the same smoke arrived today, against the shelter system that actually
exists now, and what would help most.

Two interventions present themselves, and they are not the same. The county
could add capacity---more spaces in the buildings it already operates---or it
could add capacity in better places. These sound like variations on one idea.
They are not, and the difference turns out to matter both for how many people
get indoors and for \emph{which} people do.

\runinhead{Research question.} For a fixed total shelter capacity, does the
spatial placement of that capacity change how many unsheltered residents reach a
clean-air shelter during a smoke episode, and does it change who reaches one?

I address this with an agent-based model, because the question is fundamentally
about individuals interacting with a constrained resource. A shelter that is
full turns away the next person to arrive; who arrives first depends on where
they started and how fast they walk; how fast they walk depends on their age and
their health. None of that survives aggregation. A gravity model or an
accessibility index would report an average travel distance; it would not report
that 578 spaces sat empty while 562 people were turned away, which is the
central finding here.

The model is deliberately narrow in what it claims. It measures smoke exposure
and inhaled particulate mass. It does not predict illness, hospitalization or
death, and Section~\ref{subsec:health} explains why I removed the health-risk
weighting that the project originally specified rather than repair it.

%%%%%%%%%%%%%%%%%%%%%%%%%%%%%%%%%%%%%%%%%%%%%%%%%%%%%%%%%%%%%%%%%%%%%%%%%%%
\section{Related Work and Background}
\label{sec:related}

\runinhead{Wildfire smoke and health.} Reid et al.\ \cite{reid2016critical}
review the health effects of wildfire smoke exposure and find consistent
evidence for respiratory morbidity, with weaker and less consistent evidence for
cardiovascular endpoints. DeFlorio-Barker et al.\
\cite{deflorio2019cardiopulmonary} compare wildfire and non-wildfire PM$_{2.5}$
periods among older adults across the United States. DeVries et al.\
\cite{devries2017} pool 37 studies covering approximately 1.1 million
COPD-related acute events and report a 2.5\% (95\% CI 1.6--3.4\%) increase in
COPD-related emergency visits and admissions per 10~$\mu$g~m$^{-3}$ of
PM$_{2.5}$.

That last figure is worth pausing on, because it is the kind of number that is
easy to misuse. It is an association between ambient PM$_{2.5}$ and COPD-related
events. It is \emph{not} a susceptibility multiplier comparing a person who has
COPD with a person who does not. Those are different quantities, and
Section~\ref{subsec:health} describes what happened when I tried to use
published relative risks as exposure weights.

\runinhead{Health of people experiencing homelessness.} Fazel et al.\
\cite{fazel2014health} summarize the descriptive epidemiology in high-income
countries. For respiratory prevalence specifically, the standard national
surveys are unusable: their housing-unit sampling frames structurally exclude
people who are unsheltered. Zellmer et al.\ \cite{zellmer2025} instead use
electronic health records for 20,139 adults with recent homelessness and report
diagnosed asthma at 14.9\% and COPD at 10.5\%, against 7.1\% and 3.0\% in a
housed comparison group. Brown et al.\ \cite{brown2017} study 350 homeless
adults aged 50 and over and report asthma-or-COPD at 26.3\%, essentially equal
to the all-ages figure from a large California study \cite{caspeh2023}---which
is the evidence on which I sample respiratory conditions independently of age.

\runinhead{Movement.} Bohannon and Williams Andrews \cite{bohannon2011} provide
comfortable walking speeds by age decade and sex from a meta-analysis of 41
studies and 23,111 participants. Bohannon \cite{bohannon1997} supplies the
within-population coefficient of variation. For residents with mobility
limitations, Boyce et al.\ \cite{boyce1999} report movement speeds by impairment
category from building-evacuation research. Buekers et al.\ \cite{buekers2024}
meta-analyze 25 studies comparing 1,015 people with COPD against 2,229 healthy
controls and report a usual gait-speed difference of $-19$~cm~s$^{-1}$ (95\% CI
$-28$ to $-11$), rating the evidence low quality.

I searched for an equivalent estimate for asthma and did not find one. The
literature supports lower physical \emph{activity} among adults with asthma, but
no quantitative comfortable-gait-speed decrement.
Section~\ref{subsec:results-equity} reports the resulting asymmetry as a finding
rather than concealing it.

\runinhead{Modeling platform and algorithms.} The model is built on Repast
Simphony \cite{north2013}, and uses Dijkstra's algorithm \cite{dijkstra1959}
over a graph whose edges are weighted by geodesic length computed with Karney's
algorithms \cite{karney2013}. The placement problem is a capacitated
$p$-median, which is NP-hard; I use a greedy heuristic and state in
Section~\ref{subsec:scenarios} what guarantee that does \emph{not} carry.

%%%%%%%%%%%%%%%%%%%%%%%%%%%%%%%%%%%%%%%%%%%%%%%%%%%%%%%%%%%%%%%%%%%%%%%%%%%
\section{Data and Methodologies}
\label{sec:methods}

\subsection{Data Collection}
\label{subsec:data}

Every dataset is recorded in a provenance registry in the repository with its
source, URL, retrieval date, license, SHA-256 checksum and every transformation
applied. Table~\ref{tab:data} summarizes the five inputs.

\begin{table}
\caption{Model inputs. Each row's checksum is recorded in the repository's data
registry and re-verified into every run manifest, so a result can be tied to the
exact bytes that produced it.}
\label{tab:data}
\begin{tabular}{p{2.5cm}p{3.3cm}p{1.5cm}p{3.1cm}}
\hline\noalign{\smallskip}
Dataset & Source & Retrieved & Content \\
\noalign{\smallskip}\svhline\noalign{\smallskip}
Hourly PM$_{2.5}$ & U.S.\ EPA Air Quality System, parameter 88502
  \cite{epa2020aqs} & 2026-07-24 & 4,795 rows; 312 hourly slices, no gaps \\
Street centerlines & Portland Metro RLIS & unknown & 112,070 polylines;
  89,322-node graph \\
Encampment points & City of Portland campsite reports \cite{portlandirp}
  & 2026-07-24 & 3,400 points; 2,981 used \\
Shelter inventory & Multnomah County HSD \cite{multcoshelters}; City of
  Portland \cite{portlandsrv} & 2026-07-24 & 36 facilities, 2,234 spaces \\
Population count & 2025 Tri-County Point-in-Time count \cite{pitc2025}
  & 2026-07-25 & 6,842 unsheltered \\
\noalign{\smallskip}\hline\noalign{\smallskip}
\end{tabular}
\end{table}

\subsubsection{The smoke event}

I use EPA AQS parameter 88502 rather than the Federal Reference Method series
88101 for a specific reason: I downloaded and inspected the 2020 hourly 88101
file first, and it contains \emph{no Multnomah County monitors}. Only Harney,
Klamath and Lane County sites report hourly 88101 in Oregon. Parameter 88502 is
what Oregon DEQ's Portland-area continuous monitors report and what feeds
AirNow's real-time index. The trade-off is that 88502 is designated acceptable
for index reporting but is not a reference method.

Filtering the national archive to Oregon and to the three tri-county counties
yields 4,795 hourly rows from seven monitors, two of them inside Multnomah
County. No value was altered, rounded, unit-converted, gap-filled or averaged
during extraction; the fetch script reproduces the file exactly. Averaging the
two in-county monitors by hour gives the series in Fig.~\ref{fig:event}, which
covers all 312 hours of the event window with no missing slices. EPA's own
informational qualifier \texttt{IT}, ``Wildfire -- U.S.'', is set on 1,576 rows
spanning exactly 7--19 September, which is the agency itself attesting that
these observations are smoke-influenced.

\begin{figure}[t]
  \centering
  \includegraphics[width=\textwidth]{fig1_event}
  \caption{Measured hourly PM$_{2.5}$ in Multnomah County, 7--19 September 2020,
  averaged across the two in-county regulatory monitors. The air stayed above the
  ``Unhealthy'' threshold for 194 of 312 hours and peaked at
  562.7~$\mu$g~m$^{-3}$. This is the driver of every result in the chapter; it is
  measured rather than modelled, and it is identical in all three scenarios.}
  \label{fig:event}
\end{figure}

All seven monitors use a single method---heated-inlet nephelometry---with no
method diversity and no co-located reference instrument, and the
per-observation uncertainty column is empty in all 4,795 rows. A heated inlet
volatilizes semi-volatile organic compounds, which are a large mass fraction of
fresh wood smoke, so these readings most likely \emph{understate} true
PM$_{2.5}$ during the event.

\subsubsection{Where people are}

The unsheltered population count comes from the 2025 Tri-County Point-in-Time
count \cite{pitc2025}: 10,526 people experiencing homelessness in Multnomah
County, more than 65\% of them unsheltered, giving 6,842. I use the unsheltered
subset rather than the total, because sheltered people are already indoors and
do not walk anywhere; and rather than the county's annual service count of
6,731, because that counts turnover across a year rather than people outside on
one night. The report's own caveat travels with the number: changes in how
administrative data were incorporated substantially augmented the 2025
unsheltered count, so the 2019--2025 comparison is not a clean time series and
is not presented as one.

Residents start at real locations. The City of Portland publishes campsite
reports through its Impact Reduction Program, and 3,400 of them were retrieved
through the city's open-data endpoint \cite{portlandirp}. Two limitations are
important. The feed is a rolling recent window and retains \emph{zero} records
from 2020, so 2025--26 locations serve as a spatial proxy for the 2020
distribution; the model prints this warning at the start of every run. And the
reports are complaint-driven, so they are biased toward visible camps. The
reported runs place residents at 2,981 distinct such locations.

\subsubsection{Where the shelters are}
\label{subsec:shelters}

The county's published shelter list \cite{multcoshelters} mixes five
incompatible units---beds, motel rooms, village units and pods, family units,
and unlabeled counts---so it cannot be summed directly. I converted each to
people using explicit ranges rather than point values: a bed is exactly one
person; a motel room serving ``individuals and couples'' is 1.0--1.5; a village
unit is 1.0--1.2; a family unit is 2.5--4.0 and is the weakest conversion in the
set; an unstated count is treated as people, which is the most conservative
reading because it cannot inflate the total. Applying midpoints to the full
county list gives roughly 1,900--2,200 people of year-round capacity, and the 36
modelled facilities total 2,234 spaces.

No coordinate was guessed. The raw inventory file is deliberately
coordinate-free, and every location was geocoded from a published address, with
each row recording which geocoder produced it and at what confidence. Six City
Safe Rest Village addresses were recovered from municipal announcements
\cite{portlandsrv}. Two of those publish only an intersection or a block, and
the geocoder refuses intersection queries, so they were resolved to the
corresponding hundred-block address---an approximation of a few hundred metres,
recorded as such in the confidence column rather than presented as a street
address. Three of the 36 coordinates are block- or intersection-level; the other
33 are full street addresses.

Two facilities remain excluded because neither publishes any street address:
Clinton Triangle, at 160 units the largest single site in the inventory, and
Multnomah Safe Rest Village at 28. Real capacity is therefore about 207 people
higher than modelled. Ten day centers are also excluded because none publishes a
capacity, and this exclusion cuts against my own framing: during a \emph{daytime}
smoke episode, day centers are plausibly the most relevant clean-air spaces that
exist, so scenario A understates daytime availability.

\subsection{Methods}
\label{subsec:methods}

\subsubsection{Platform, languages and tools}

The simulation is written in Java~17 on Repast Simphony 2.11 \cite{north2013},
built with Gradle, and uses the GeographicLib library for geodesic computation
and the JTS Topology Suite for spatial indexing. Analysis, figure generation,
scenario construction and the placement optimizer are Python 3.14 using pandas
and Matplotlib. All source, data, configuration and run manifests are in the
repository named in the code-availability note.

\runinhead{Use of AI.} I directed this research, made all research decisions,
and wrote this manuscript. I used Claude (Anthropic) as a coding and
verification assistant: it helped write Java and Python source, build the
data-acquisition and analysis scripts, draft documentation, and check
bibliographic records against publisher pages. The citation audit reported in
Section~\ref{subsec:health}---which found that two parameters in my original
design were attributed to papers that do not report them, and one to a paper
that does not exist---was carried out with that assistance and then verified by
me against the publisher records. I reviewed, revised and approved all outputs
and take full responsibility for the final text and for every number in it.

\subsubsection{Agents}

Each resident is an agent with a state, a position, an origin, a walking speed
and a set of sampled attributes. The state set is $\{$\textsc{pre-evac},
\textsc{en-route}, \textsc{sheltered}, \textsc{unreachable},
\textsc{refused-all-full}$\}$. A resident begins in \textsc{pre-evac} at their
encampment and departs only when both conditions in Eq.~\ref{eq:evac} hold:

\begin{equation}
\label{eq:evac}
C(t) \geq 55.5\ \mu\mathrm{g\ m}^{-3}
\quad\wedge\quad
\exists\, s \in S : \mathrm{open}(s, t)
\end{equation}

\noindent where $C(t)$ is the county PM$_{2.5}$ concentration at time $t$ and
$S$ is the shelter set. The threshold is the lower bound of EPA's ``Unhealthy''
category. I chose that particular breakpoint because it is stable across both
the pre-2024 and post-2024 breakpoint tables, unlike the categories above it.
Because index categories are defined on 24-hour averages while this model counts
hourly observations, the derived metric is named and defined as hours above a
concentration and is never described as an index category.

Attributes are drawn from the sources in Table~\ref{tab:params}. One engineering
detail is load-bearing. Repast shuffles agent execution order each tick using
its default random stream, which is the same stream that assigns each resident
an encampment. Drawing a single additional number from it would silently
regenerate the entire population. The sampler therefore uses a private stream
derived from the run seed as $\mathrm{seed} \times 1000003 + 17$; the verified
consequence is that adding the whole attribute layer left the archived baseline
run byte-identical.

\begin{table}
\caption{Sampled resident attributes with their sources and realized marginals
at seed 42. Every value is either measured (M) or from published literature (L);
the model refuses to start if a variable in either class lacks a resolvable
identifier for its source.}
\label{tab:params}
\begin{tabular}{p{3.0cm}p{1.9cm}p{1.5cm}p{3.5cm}}
\hline\noalign{\smallskip}
Attribute & Target & Realized & Source \\
\noalign{\smallskip}\svhline\noalign{\smallskip}
Age 18--44 / 45--64 / 65+ (\%) & 52.7 / 42.3 / 5.0 & 52.8 / 42.0 / 5.2 &
  Pathways Study 2026, $N=541$, local (M) \\
Sex M / F / other (\%) & 68.4 / 29.3 / 2.3 & 68.6 / 29.2 / 2.2 &
  2019 Multnomah PIT (M) \\
Mobility limitation (\%) & 19.2 & 19.9 & 2019 PIT; a lower bound (M) \\
Asthma (\%) & 15.0 & 14.8 & Zellmer et al.\ \cite{zellmer2025} (L) \\
COPD (\%) & 10.5 & 10.8 & Zellmer et al.\ \cite{zellmer2025} (L) \\
Chronic physical condition (\%) & 39.1 & 39.6 & Pathways 2026, local (M) \\
Walking speed (m~s$^{-1}$) & --- & 1.28 mean & Bohannon and Williams Andrews
  \cite{bohannon2011} (L) \\
\noalign{\smallskip}\hline\noalign{\smallskip}
\end{tabular}
\end{table}

\subsubsection{Walking speed}

Walking speed is the only attribute permitted to change an outcome; everything
else is a reporting stratum. A resident without a mobility limitation draws from
a truncated normal centred on the Bohannon age-by-sex cell for their decade and
sex, with standard deviation $\mathrm{CV} \times \mu$ where $\mathrm{CV}=0.13$.
A resident with COPD has that mean shifted additively by $-0.19$~m~s$^{-1}$,
matching how Buekers et al.\ report the difference. A resident with a mobility
limitation instead draws from $\mathcal{N}(0.95, 0.32)$, \emph{replacing} the
age-by-sex mean rather than scaling it, and the COPD decrement is not stacked on
top of it, because the impaired categories already embed a slower walker:

\begin{equation}
\label{eq:speed}
v_i \sim
\begin{cases}
\mathcal{N}_{[0.4,\,2.2]}(0.95,\ 0.32) & \text{mobility limited}\\[3pt]
\mathcal{N}_{[0.4,\,2.2]}\!\left(\mu_{a,s} - 0.19\,\delta^{\mathrm{COPD}}_i,\
  0.13\,\mu_{a,s}\right) & \text{otherwise}
\end{cases}
\end{equation}

\noindent where $\mu_{a,s}$ is the published mean for age decade $a$ and sex $s$,
and $\delta^{\mathrm{COPD}}_i \in \{0,1\}$. The standard deviation comes from
Bohannon's \emph{within-population} coefficient of variation, never from the
2011 meta-analysis's confidence intervals. Those are between-study intervals;
using them would understate the spread of individual walking speeds by a factor
of three to five. Fig.~\ref{fig:speed} shows the resulting distributions.

\begin{figure}[t]
  \centering
  \includegraphics[width=\textwidth]{fig6_speed}
  \caption{Sampled comfortable walking speeds by group. The mobility-limited
  distribution is a replacement for, not a rescaling of, the age-by-sex
  distribution. Asthma appears in no group here, because no published gait-speed
  decrement exists for it; that absence propagates directly into the equity
  result in Fig.~\ref{fig:equity}.}
  \label{fig:speed}
\end{figure}

\subsubsection{Movement and routing}

The street layer is reprojected to WGS84 at load and converted into an
undirected graph of 89,322 nodes and 112,070 edges, with each edge weighted by
its geodesic polyline length \cite{karney2013}. The shapefile's own length
attribute is deliberately not trusted, because its unit is undocumented. One
Dijkstra tree \cite{dijkstra1959} is computed per shelter at setup; because the
graph is undirected, the distance from any agent's node to a shelter is then a
constant-time lookup, which is what allows a 312-hour run at $n=6{,}842$ to
complete in under a minute. Each tick, an agent consumes a budget of
$v_i \times 60 \times \Delta t$ metres along its routed polyline, finishing any
partial segment with a geodesic direct solve, so movement follows the street
network at exact arc length rather than interpolating in a straight line.

\runinhead{A defect in the street data.} During validation I found that 27 node
identifiers were claimed by different street features at locations 9 to 18.5~km
apart. The shapefile contains no multi-part geometry, so this was a pure
attribute defect. Its effect was severe: the affected edges were \emph{weighted}
in metres but \emph{spanned} kilometres, so the shortest-path search
under-costed them, agents routed through them, and then traversed the gap in a
straight line. Twenty-two of fifty agents in a test run had inflated journeys,
with travel distance up to 15 times too large and inhaled dose three times too
large, and shelter \emph{choice} could be wrong rather than merely shelter
distance.

I corrected it without deleting anything. Every duplicate claim beyond the first
is either reattached to a geometrically coincident junction within 10~m (4
cases) or split onto a new synthetic node at its true location (23 cases), with
every correction written into the run manifest. The 100~m clustering threshold
sits in an empty band in the data: legitimate endpoint scatter is sub-metre and
the smallest corrupt displacement is about 1.65~km, so any threshold between
those regimes yields the identical graph. After correction, impossible edges
fell from 50 to 0, the maximum endpoint gap from 18.5~km to 11.9~m, and the
component structure was unchanged.

\subsubsection{Exposure, dose and health risk}
\label{subsec:health}

The model computes three quantities and never multiplies them together by
accident. Exposure is the environmental integral, a property of the air:

\begin{equation}
\label{eq:exposure}
E_i = \sum_{t} C(t)\,\Delta t
\qquad [\mu\mathrm{g\ m}^{-3}\,\mathrm{h}]
\end{equation}

Inhaled dose is a property of the person, differing from exposure only by
ventilation rate, which depends on activity and on nothing else:

\begin{equation}
\label{eq:dose}
D_i = \sum_{t} C(t)\,\mathrm{IR}\!\left(\alpha_i(t)\right)\Delta t,
\qquad
\mathrm{IR} =
\begin{cases}
1.62\ \mathrm{m^3\,h^{-1}} & \alpha_i(t) = \text{walking}\\
0.61\ \mathrm{m^3\,h^{-1}} & \text{otherwise}
\end{cases}
\end{equation}

\noindent with rates from the EPA Exposure Factors Handbook \cite{epaefh2011}.
Health risk is a third, separate quantity, $R_i = D_i \cdot w_i$, and
$w_i = 1.0$ for every resident by design. Accrual stops when a resident is
admitted; residents who are walking, stranded or turned away are all still
outdoors and keep accruing.

The reason $w_i$ is inert deserves a full account, because it is the single
largest methodological change I made. The project originally specified an
exposure weighted by relative risks: 1.45 for adults 65 and over, attributed to
Di et al.\ \cite{di2017}, and 1.80 for chronic obstructive pulmonary disease,
attributed to ``Anderson et al.\ 2013''. Checking those against the publisher
records found three problems.

First, Di et al.\ is a real and excellent study of 60.9 million Medicare
beneficiaries, and it reports a hazard ratio of 1.073 per 10~$\mu$g~m$^{-3}$ of
annual PM$_{2.5}$---not 1.45. More fundamentally, its cohort is \emph{entirely}
aged 65 and over, so the study design cannot produce an age contrast at all.
Second, no paper matching ``Anderson et al.\ 2013'' with that finding exists.
The nearest record is a 2013 study whose first author is Atkinson
\cite{atkinson2013}, published in a different journal, examining cardiovascular
disease incidence, and reporting no COPD effect-modification estimate. Third,
and independent of the citations, multiplying an exposure by a relative risk is
a category error: a relative risk is a ratio of outcome \emph{rates} between
groups, not a multiplier on \emph{exposure}. The correct object would be a ratio
of exposure--response coefficients, which is a different quantity, and the same
underlying data yield an age multiplier anywhere between 1.008 and 2.5 depending
only on which scale is chosen.

There is also direct evidence against the specific value. A meta-analysis of
eight studies reports an elderly-to-adult relative-risk ratio of 1.008 (95\% CI
0.996--1.020)---a null result. The proposed 1.45 was not merely unsourced; the
best available evidence suggests the true value is close to one.

I therefore set all susceptibility weights to 1.0 and made \emph{stratified
reporting} the primary equity metric: outcomes are reported separately for each
group rather than collapsed into a weighted index that no source can justify.
The weighting slot remains in the code so that a sourced coefficient has exactly
one place to land, and so a reader can see that risk weighting is switched off
rather than absent. For the same reason, inhalation rate is held constant across
health status: no defensible population-specific ventilation multiplier for
adults with asthma or COPD during walking was identified, and any figure would
have been an assumption wearing a citation.

\subsubsection{The three scenarios}
\label{subsec:scenarios}

The scenarios are not three guesses. Each answers the question the previous one
measured. Scenario A is a \emph{measurement} rather than a treatment: its only
job is to reveal which constraint binds. It reported capacity, so B relieves
capacity and changes nothing else. B then revealed a second constraint, so C
spends the identical total capacity differently.

\begin{itemize}
\item \textbf{A --- today.} All 36 existing shelter facilities at their real
geocoded coordinates with their real capacities: 2,234 spaces.
\item \textbf{B --- capacity meets demand.} Every real facility scaled by
$6842/2234 = 3.0627$, apportioned by largest remainder so the system total is
exactly 6,842. Coordinates, facility count and relative facility size unchanged.
\item \textbf{C --- the same capacity, better placed.} Real facilities grow by
only $1.5\times$; the remaining capacity becomes ten new facilities at
algorithmically chosen street-network nodes. Total capacity is held identical to
B at 6,842.
\end{itemize}

Because B and C hold total capacity exactly equal, a B$\rightarrow$C difference
isolates \emph{where} the marginal capacity sits and nothing else. C never moves
an existing facility: all 36 stay at their real coordinates, because a real
shelter system cannot be picked up and set down elsewhere.

The ten new locations come from a capacitated greedy $p$-median over candidate
street nodes thinned to one per 600~m grid cell. The objective minimizes
$\sum_{i,j} d_{ij} q_{ij}$, the network distance times residents served, plus a
declared penalty per unfilled space. Facilities are placed largest-first and
served demand is decremented before the next placement. This is a heuristic: it
does not carry the $(1-1/e)$ submodularity guarantee, because shelter catchments
overlap. The result is a good placement, not a provably near-optimal one, and
the $1.5\times$ factor and the choice of ten sites are policy parameters rather
than measured quantities.

\subsubsection{Experimental design}

Everything except the shelter file is held identical across arms: the same
6,842 residents with the same identifiers, start coordinates, ages, sexes,
mobility statuses, respiratory conditions and walking speeds; the same smoke
field; the same street graph; the same parameters. Each arm was run with nine
random seeds (42--50), giving 27 runs. Every reported result is seed 42, with
the range across all nine seeds given alongside.

%%%%%%%%%%%%%%%%%%%%%%%%%%%%%%%%%%%%%%%%%%%%%%%%%%%%%%%%%%%%%%%%%%%%%%%%%%%
\section{Results}
\label{sec:results}

\subsection{Verification}
\label{subsec:verification}

Before reporting outcomes I report what was checked, including several tests
whose answers can be worked out by hand.

The exposure integral was recomputed in Python directly from the raw EPA file,
independently of the model. A resident who never reaches shelter should
accumulate the full-window integral of 54,002.7~$\mu$g~m$^{-3}$~h; the model
reports 54,002.8, a ratio of 1.0000. The same residents show a mean of
173.09~$\mu$g~m$^{-3}$, a peak of 562.7, and 194.0 hours above threshold, each
matching the independently computed value exactly.

Routing was validated against a separate Python implementation of Dijkstra's
algorithm over the same graph, which reproduces the Java network distances
exactly across five tests. Realized walking speeds recovered from exported
distance and time fall in 1.300--1.376~m~s$^{-1}$, inside the published Bohannon
bounds. Re-running at seed 42 reproduces output bit-for-bit.

A specific defect motivated one further check. A resident turned away at a full
shelter originally re-planned its route from its immutable start node, which
meant it walked back to its encampment before setting out again---inflating both
distance and dose. The bug was invisible at the 50-agent test scale because no
shelter ever filled there. After the fix, every run is checked against
$\mathrm{walked} \leq \mathrm{planned} + \mathrm{snap\ gap} + 200$~m for every
agent; in a capacity-binding reference run where 250 of 400 agents were refused
at least once, the maximum unexplained distance was 8.9~m.

Across the 27 production runs, six cross-run invariants hold: manifest seed and
scenario code match the run; the working tree was clean for every run; the data
version is identical across all nine seeds within each arm; the source-file
checksum set is identical across all 27; each run exports exactly 6,842 resident
records; and within each seed the population is byte-identical across arms,
verified by SHA-256 over the joined attribute vector.

\subsection{Access}
\label{subsec:results-access}

Table~\ref{tab:results} gives the headline outcomes and
Fig.~\ref{fig:outcomes}(a) shows the access result.

\begin{table}
\caption{Outcomes at seed 42, with the range across all nine seeds in brackets.
No range overlaps between arms on any metric. B and C hold identical total
capacity, so the B$\rightarrow$C difference is attributable to placement alone.}
\label{tab:results}
\begin{tabular}{p{3.9cm}p{2.3cm}p{2.3cm}p{2.3cm}}
\hline\noalign{\smallskip}
 & A: today & B: more beds & C: better placed \\
\noalign{\smallskip}\svhline\noalign{\smallskip}
Facilities & 36 & 36 & 46 \\
Total spaces & 2,234 & 6,842 & 6,842 \\
Sheltered & 2,060 (30.1\%) & 6,264 (91.6\%) & \textbf{6,570 (96.0\%)} \\
\quad range, 9 seeds & [2,053--2,064] & [6,257--6,268] & [6,563--6,574] \\
Turned away, all full & 4,766 & 562 & \textbf{256} \\
Spaces left empty & 174 & \textbf{578} & 272 \\
No shelter reachable & 16 & 16 & 16 \\
Mean distance walked (m) & 18,260 & 7,938 & \textbf{5,689} \\
Mean hours above threshold & 135.8 & 17.5 & \textbf{8.6} \\
Mean inhaled dose ($\mu$g) & 23,374 & 3,056 & \textbf{1,534} \\
Person-hours above threshold & 928,934 & 119,921 & \textbf{59,060} \\
\noalign{\smallskip}\hline\noalign{\smallskip}
\end{tabular}
\end{table}

Today's system shelters 2,060 of 6,842 residents, or 30.1\%. Capacity binds
hard: 4,766 residents are turned away because every reachable facility is full,
and 33 of the 36 facilities finish at capacity. Raising capacity to meet demand
at the same locations lifts this to 6,264 (91.6\%), a factor of 3.04. Total
exposure falls 87.3\% and mean walking distance falls 56.5\%.

\begin{figure}[t]
  \centering
  \includegraphics[width=\textwidth]{fig2_outcomes}
  \caption{(a) Residents reaching a shelter. (b) Spare capacity and unmet need in
  the same run. Panel (b) is the central result: scenario B leaves 578 spaces
  empty while turning 562 people away, so B has no shortage---it has exactly one
  space per person---and fails anyway, because the spaces are where the buildings
  are rather than where the people are.}
  \label{fig:outcomes}
\end{figure}

\subsection{The geography failure}
\label{subsec:results-geo}

Scenario B leaves \textbf{578 spaces empty while turning 562 people away}. Those
two numbers are nearly equal, and they are the reason scenario C exists. B has
enough capacity for the entire population and still fails 562 people, because
the added capacity went to the buildings that already exist rather than to the
places where people are. That is a geography failure, and the design separates it
cleanly from a capacity failure, because capacity is not binding in aggregate in
B.

Scenario C spends the identical 6,842 spaces. Refusals fall from 562 to 256,
empty spaces from 578 to 272, mean walking distance by 28.3\%, total exposure by
a further 50.7\%, and mean inhaled dose from 3,056 to 1,534~$\mu$g.
Fig.~\ref{fig:map} shows the spatial arrangement that produces this.

\begin{figure}[t]
  \centering
  \includegraphics[width=\textwidth]{fig4_map}
  \caption{Resident start points (grey) against facility locations, sized by
  capacity. In (a) the existing facilities cluster away from much of the demand;
  in (b) the ten new sites are placed into the demand surface itself. The
  existing 36 facilities are unmoved between panels---only the new capacity is
  placed.}
  \label{fig:map}
\end{figure}

Relative to today, scenario C shelters 3.19 times as many residents, reduces
total exposure by 93.8\%, person-hours above threshold by 93.6\%, and mean
distance walked by 68.8\%. The ranking of interventions is therefore
unambiguous: adding capacity is the first-order effect and placing it well is
the second-order effect---but the second-order effect costs no additional
spaces.

\subsection{Equity}
\label{subsec:results-equity}

The second finding is one that an aggregate analysis cannot produce, because an
evaluation tracking only total access would score scenario B as a near-complete
success. Table~\ref{tab:equity} and Fig.~\ref{fig:equity} report access by
group.

\begin{table}
\caption{Percentage of each group reaching a shelter, seed 42. Adding capacity at
existing sites widens the mobility gap from 13.0 to 24.5 percentage points;
spending the same capacity on well-placed sites returns it to 12.9 while lifting
the slowest group from 71.9\% to 85.7\%.}
\label{tab:equity}
\begin{tabular}{p{4.3cm}p{1.5cm}p{1.5cm}p{1.5cm}p{1.5cm}}
\hline\noalign{\smallskip}
Group & Share (\%) & A (\%) & B (\%) & C (\%) \\
\noalign{\smallskip}\svhline\noalign{\smallskip}
Everyone & 100.0 & 30.1 & 91.6 & 96.0 \\
Walks without difficulty & 80.1 & 32.7 & 96.4 & 98.6 \\
Mobility limitation & 19.9 & 19.7 & 71.9 & 85.7 \\
Age 65 and over & 5.2 & 22.4 & 82.4 & 89.8 \\
COPD & 10.8 & 22.2 & 86.2 & 93.8 \\
Asthma & 14.8 & 29.2 & 90.6 & 95.7 \\
Chronic physical condition & 39.6 & 30.2 & 91.1 & 95.8 \\
\noalign{\smallskip}\svhline\noalign{\smallskip}
Mobility gap (points) & --- & 13.0 & 24.5 & 12.9 \\
\noalign{\smallskip}\hline\noalign{\smallskip}
\end{tabular}
\end{table}

\begin{figure}[t]
  \centering
  \includegraphics[width=\textwidth]{fig3_equity}
  \caption{Access by mobility status. Pouring 4,608 additional spaces into the
  same buildings widens the gap, because extra capacity at an existing site is
  captured first by whoever can walk there fastest; well-placed capacity narrows
  it again. The gap is narrowed, not closed---14.3\% of residents with mobility
  limitations are still outside in scenario C.}
  \label{fig:equity}
\end{figure}

Capacity expansion alone \emph{widens} the equity gap. In scenario A the
difference in access between residents who walk without difficulty and residents
with a mobility limitation is 13.0 percentage points. Adding 4,608 spaces to the
same buildings takes it to 24.5, because additional capacity at an existing site
is claimed first by whoever reaches it fastest. Spending that identical capacity
on well-placed new sites brings the gap back to 12.9 while raising the slowest
group from 71.9\% to 85.7\%. The same pattern holds for residents aged 65 and
over (22.4, 82.4, 89.8) and for residents with COPD (22.2, 86.2, 93.8).

Asthma shows almost no access penalty while COPD shows a large one. This
asymmetry is a direct consequence of the evidence, not an omission. A diagnosis
can only affect an outcome in this model by affecting movement, because
diagnosis is never a dose multiplier; COPD carries a published gait-speed
decrement and asthma does not. Had I borrowed the COPD estimate for asthma to
make the treatment symmetric, I would have manufactured a finding. The asymmetry
is therefore evidence that the model is not inventing effects.

\subsection{Robustness}
\label{subsec:results-robust}

Fig.~\ref{fig:seeds} shows every one of the 27 runs. The largest seed-to-seed
spread on residents sheltered is 11 people; the smallest gap between two arms is
306. The between-arm signal is approximately 28 times the within-arm noise, and
no range overlaps between arms on any headline metric.

\begin{figure}[t]
  \centering
  \includegraphics[width=\textwidth]{fig5_seeds}
  \caption{Residents sheltered in all 27 runs, on a logarithmic scale. The
  vertical separation between arms is roughly 28 times the seed-to-seed spread
  within an arm, so the ordering A $<$ B $<$ C is not an artifact of any
  particular random draw.}
  \label{fig:seeds}
\end{figure}

One detail falls out of the design as a free consistency check. The count of
residents who can reach no shelter at all is identical across all three arms
within each seed---16, 16, 16 at seed 42. Those residents sit on street-graph
components containing no facility, so no amount of capacity or re-placement
reaches them, and the number varies only with which residents were sampled.

\subsection{Calibration}
\label{subsec:calibration}

Against the one observed occupancy record I could locate---a contemporaneous
news report of approximately 130 of 198 spaces occupied on 16 September
2020---a historical reference configuration of the model fills all 198. That is
a 1.52 times over-prediction. I attribute it to the assumption that every
resident knows the shelters exist, against local survey evidence that 65\% of
unsheltered respondents had never heard of them. Every ``reached a shelter''
figure in this chapter is therefore an \emph{upper bound}, and the wording used
throughout is that placement improves outcomes \emph{under the modelled
assumptions}, never that the model reproduces what happened.

%%%%%%%%%%%%%%%%%%%%%%%%%%%%%%%%%%%%%%%%%%%%%%%%%%%%%%%%%%%%%%%%%%%%%%%%%%%
\section{Conclusion}
\label{sec:conclusion}

\subsection{Applications}

The practical implication is a ranking. If Multnomah County can add clean-air
shelter capacity, that dominates every other intervention modelled here: it
moves access from 30.1\% to 91.6\% and cuts total exposure by 87.3\%. But the
second finding is the one with no budget attached. Once capacity is adequate in
aggregate, \emph{where} the marginal capacity sits determines whether it is
usable, and the model puts a number on that: the same 6,842 spaces halve
refusals and halve inhaled dose depending only on placement.

The equity result sharpens this into an operational warning. A capacity
expansion evaluated on aggregate access would look like a success while the gap
between residents who walk easily and residents who do not nearly doubled. Any
real expansion should be evaluated on stratified access, not on totals.

The method generalizes beyond smoke. The same structure---a hazard field, a
population with heterogeneous mobility at real locations, and a capacitated set
of refuges on a real network---applies to extreme-heat cooling centers, warming
shelters, and emergency distribution points. The specific finding that
capacity-at-existing-sites widens mobility gaps is a hypothesis those settings
could test.

\subsection{Concerns and Next Steps}

\runinhead{What this does not claim.} No health outcome is modelled. The model
measures exposure and inhaled particulate mass, and predicts no illness,
hospitalization or death. The ten new sites are street-network nodes, not
buildings: no zoning, cost, staffing or indoor-air-filtration analysis stands
behind them. The $1.5\times$ expansion factor, the choice of ten sites, and B's
uniform $3.06\times$ scale-up are policy parameters, not measurements; the
comparison establishes direction, and magnitudes would move with different
values.

\runinhead{Limitations that cut against the result.} Several known problems work
against the conclusions rather than for them, and are stated for that reason.
Ten day centers are excluded because none publishes a capacity, yet in a daytime
smoke episode they are plausibly the most relevant clean-air spaces that
exist---so scenario A understates present-day daytime availability, which is the
baseline the other arms are compared against. Every mobility-limited resident is
assigned the \emph{fastest} of the published impaired-movement categories, so the
modelled mobility penalty is conservative and the real access disparity is
probably larger than reported. Two real facilities totalling about 207 people of
capacity are missing entirely because neither publishes an address.

\runinhead{Limitations of the inputs.} The PM$_{2.5}$ field is spatially uniform,
because with two in-county monitors any interpolation would manufacture gradients
the data cannot support; a consequence is that placement can only help by
reducing time outdoors, never by moving people into cleaner air. The monitors are
non-reference instruments that most likely understate true concentrations during
fresh wood smoke. Encampment locations are 2025--26 complaint-driven reports used
as a spatial proxy for 2020. Sex and mobility distributions are 2019 inside an
otherwise 2026 study, and asthma and COPD prevalences are imported from a
Minnesota cohort because no local equivalent exists. Admission is
order-dependent: residents are served in shuffle order rather than by need, which
matters most in scenario A where capacity binds hardest.

\runinhead{Next steps.} The most valuable single addition is an awareness model.
The one sourced behavioral parameter available---that 65\% of surveyed
unsheltered residents had never heard of the clean-air shelters---is not yet
implemented, and implementing it as a Bernoulli draw with a swept rate would
convert every access figure from an upper bound into an interval. Second,
order-independent two-phase admission would remove the shuffle-order dependence
that currently affects scenario A. Third, a needs-based triage policy could be
compared against the current first-come-first-served rule, which would let the
equity result be attacked by policy rather than only by geometry. Finally, the
$1.5\times$ and ten-site parameters should be swept to trace the full trade-off
curve between enlarging existing facilities and building new ones.

\begin{acknowledgement}
This work was completed through the National Science Foundation Research
Experiences for Undergraduates program, Computational Modeling Serving Portland,
hosted by Portland State University. I thank Prof.\ Christof Teuscher for his
mentorship and guidance throughout the program.

I directed the research, made all research decisions, and wrote the manuscript.
Claude (Anthropic) assisted with coding, data-acquisition and analysis
scripting, verification tooling, and documentation drafting, and assisted with
the citation audit reported in Section~\ref{subsec:health}. I reviewed, revised
and approved all outputs and take full responsibility for the final text.
\end{acknowledgement}

\runinhead{Code and data availability.} The full simulation source, analysis
scripts, input data, configuration files and all 27 run manifests are available
at \url{https://github.com/fxa28196/REU} under the MIT License. Every result in
this chapter was produced from a clean working tree; each run manifest records
its git commit, random seed, Java and Repast versions, every parameter value,
and a SHA-256 checksum for every input file. The reported results use seeds 42
through 50, with seed 42 as the reported run. The EPA air-quality data and the
City of Portland campsite reports are public and their retrieval scripts are
included; the street centerline file is redistributed as received, and its
original release version and license could not be recovered, so the repository's
licence file states its terms as unverified rather than relicensing it.

%%%%%%%%%%%%%%%%%%%%%%%%%%%%%%%%%%%%%%%%%%%%%%%%%%%%%%%%%%%%%%%%%%%%%%%%%%%
\bibliography{references}
\bibliographystyle{plain}

\end{document}
